% =============================================================================
% D8 — Financement et cycles d'infrastructure — v3 (mai 2026)
% Réécriture : ajout table + « matériau révèle » + sources contemporaines
% + nouvelles refs (Esty, Dell'Ariccia, Falato, Lerner-Nanda, Janeway, Flamm)
% + scrap iron clause + progression observation→théorie→hypothèse
% Hypothèses H-D8.1/2/3 INCHANGÉES (cohérentes avec tableau récapitulatif)
% =============================================================================

\clearpage
\subsection{Financement et cycles d'infrastructure}\label{sec:d8-finance}

La capacité technique d'une infrastructure informationnelle ne suffit pas à son déploiement. L'analyse historique montre que le déploiement est souvent conditionné par l'existence d'un véhicule de financement adapté à la structure temporelle des coûts et des revenus. Cette dimension examine les mécanismes par lesquels la structure de financement détermine le rythme et la forme du déploiement de l'infrastructure informationnelle.

\subsubsection{Le champ~: la structure de financement comme déterminant du déploiement}

Le point de départ est un constat empirique que l'analyse historique a documenté dans cinq cas distincts. Le câble transatlantique de Cyrus Field était techniquement réalisable dès~1857, mais il fallut cinq tentatives, deux faillites et neuf ans pour qu'un câble fonctionnel soit posé en~1866 \parencite{gordon_thread_2002}. La machine analytique de Babbage était, selon la reconstruction de \textcite{swade_difference_2001}, techniquement constructible avec la mécanique de précision disponible dans les années~1840~; elle ne fut jamais achevée, et le gouvernement britannique ne la refinança pas. La tabulatrice d'Hollerith fonctionnait parfaitement dès~1890~; elle ne devint une industrie que lorsqu'IBM imposa le modèle de location avec consommable exclusif \parencite{cortada_before_1993}. Le téléphone de Bell créait un marché mesurable dès~1877~; il ne devint une infrastructure nationale qu'après la reconsolidation Morgan-Vail de~1907--1913. Le réseau électrique d'Insull produisait un kilowattheure moins cher que les dynamos privées dès les années~1890~; il ne se généralisa qu'avec l'invention du cadre réglementaire des \emph{utilities}, qui garantissait à l'investisseur un rendement sur capital en échange d'une obligation de service \parencite{hughes_networks_1983}. À Londres, l'\emph{Electric Lighting Act} de~1882 gela l'investissement privé pendant vingt ans par une clause de rachat municipal à la valeur de la ferraille (\emph{scrap iron clause})~: soixante-neuf autorisations provisoires en~1883, zéro en~1885 \parencite{hughes_networks_1983}. Dans chaque cas, la capacité technique précédait le déploiement~; dans chaque cas, c'est un arrangement financier --- pas une innovation technique --- qui débloquait ou bloquait la diffusion.

\markNEO Le benchmark néoclassique suppose que la finance ne compte pas. Le théorème de \textcite{modigliani_cost_1958} établit que, sous des conditions idéales --- pas de taxes, pas de coûts de faillite, pas d'asymétrie informationnelle, marchés complets ---, la structure de financement n'affecte pas la valeur de la firme. \NEO{La proposition d'irrélevance de la structure financière est le point de départ logique~: si elle tenait, D8 n'existerait pas comme dimension analytique, et le déploiement dépendrait exclusivement des propriétés de la technologie et de la demande.} Or les six cas qui précèdent montrent que D8 existe~: le déploiement n'est pas réductible à D1--D7. \textcite{falato_intangible_2022} en formalisent l'une des raisons pour le cas contemporain~: les firmes à actifs intangibles accumulent du cash parce que ces actifs ne peuvent pas servir de collatéral~--- la structure financière de la firme est endogène à la nature de son capital. Le lien D3~$\to$~D8 est direct~: la forme du capital détermine les conditions de son financement.

\markEVO \textcite{perez_technological_2002} propose le cadre le plus abouti pour penser le rapport entre finance et déploiement d'infrastructure. \EVO{Chaque révolution technologique traverse quatre phases~--- irruption, frénésie, synergie, maturité~--- et la transition de la frénésie à la synergie passe par une crise financière qui n'est pas un accident mais un mécanisme de transfert~: le capital spéculatif finance l'installation de l'infrastructure, la crise en transfère la propriété à des opérateurs qui la déploient.} Le pattern est documenté pour les canaux (1790s), les chemins de fer (1840s), l'acier et l'électricité (1890s), et les dotcoms (2000). \textcite{janeway_doing_2012}, praticien du capital-risque pendant quarante ans, confirme Perez par un registre différent~: l'investissement spéculatif est une condition nécessaire pour construire l'infrastructure excédentaire qui servira de base à la vague suivante. \textcite{reinhart_this_2009} le confirment par un troisième registre~: le syndrome \og \emph{this time is different}\fg{}, documenté sur huit siècles de crises financières, accompagne systématiquement le déploiement des nouvelles infrastructures.

Le champ de D8 se formule alors ainsi~: quels mécanismes causaux lient la structure de financement au rythme et à la forme du déploiement de l'infrastructure informationnelle~? La question n'est pas descriptive (quels instruments financiers sont utilisés~?) mais analytique (pourquoi certains véhicules financiers permettent le déploiement et d'autres non~?). D7 traite la forme organisationnelle une fois le capital levé~; D8 traite les conditions de la levée elle-même.

\paragraph{Contributions propres de la section}
Cette section distingue trois mécanismes que la littérature traite généralement comme un seul. H-D8.1 reprend \textcite{perez_technological_2002} sur l'investissement spéculatif comme étape fonctionnelle du déploiement, en explicitant ses conditions d'activation. H-D8.2 traite l'adéquation du véhicule financier à la structure temporelle des coûts comme variable critique, avec le cas Babbage comme anti-pattern méthodologique~--- la machine analytique échoue non par défaut technique mais par absence d'un véhicule de recapitalisation. H-D8.3 identifie une sélection adverse entre infrastructures concurrentes selon la finançabilité, dont le marqueur contemporain est le basculement \emph{venture capital} vers \emph{project finance}.


\begin{table}[!ht]
\caption{Cas historiques au prisme de trois propriétés du véhicule financier.}\label{tab:d8-vehicules}
\centering\footnotesize
\renewcommand{\arraystretch}{1.35}
\begin{tabularx}{\textwidth}{@{}p{2.0cm}p{2.8cm}p{2.0cm}p{2.0cm}p{2.0cm}X@{}}
\toprule
\textbf{Cas} & \textbf{Véhicule} & \textbf{L'actif est-il collatéralisable~?} & \textbf{Le flux de revenus est-il identifiable~?} & \textbf{Horizon coûts / revenus} & \textbf{Résultat} \\
\midrule
Câble 1857--66 & Concession & Oui (câble physique) & Oui (arbitrage intercontinental) & Ponctuel / récurrent & Faillite puis succès après recapitalisation \\
Babbage 1823--42 & Subvention publique unique & Oui (machine) & Non (pas de marché) & Cumulatif / nul & Abandon \\
IBM 1920s--40s & Location + consommable exclusif & Oui (tabulatrice) & Oui (cartes, marge 60--80\,\%) & Capex converti en opex récurrent & Dividendes pendant la Dépression \\
Bell 1876--94 & Brevet de principe & Non (droit incorporel) & Oui (rente légale, 46\,\%) & 18~ans de monopole & Stagnation puis concurrence \\
Insull 1890s & Régulation \emph{utility} & Oui (centrale + réseau) & Oui (kWh, rendement garanti) & Irréversible $\to$ quasi-bond & Déploiement \\
Londres 1882 & Concession + \emph{scrap iron clause} & Oui mais valeur amputée & Oui mais horizon tronqué (21~ans) & Amputé au rachat & Gel de l'investissement \\
Transpacifique 1902 & Concession & Oui & Non (pas de corridor commercial) & Ponctuel / absent & Échec commercial \\
\bottomrule
\end{tabularx}
\renewcommand{\arraystretch}{1.0}
\end{table}


\subsubsection{Ce que le matériau révèle}

Avant de formuler les mécanismes, il est utile de poser ce que l'analyse historique fait apparaître quand on la relit à travers la grille du tableau~\ref{tab:d8-vehicules}.

Un premier constat concerne la récurrence du cycle. La séquence investissement massif, crise, recapitalisation apparaît dans le câble transatlantique (1857--1866), le télégraphe américain (vingt compagnies dans les années~1850, consolidation Western Union en~1866), et le téléphone (monopole de brevet, puis concurrence explosive après~1894, puis reconsolidation Morgan-Vail en~1907). Dans chaque cas, l'infrastructure survit à la crise~; ce sont les financeurs initiaux qui sont éliminés. Le pattern est absent là où les rendements de réseau manquent~: le marché des calculatrices mécaniques, sans rendements croissants ni coûts fixes prohibitifs, n'a connu ni bulle ni consolidation.

Un deuxième constat concerne la variable discriminante. La confrontation des cas fait apparaître que ce qui sépare le succès de l'échec n'est pas la qualité technique mais la présence d'un flux de revenus identifiable. Babbage et le câble ont des techniques comparables en qualité~; le câble a un flux (l'arbitrage intercontinental) et Babbage n'en a pas. IBM survit à la Dépression par les cartes (un flux récurrent). Londres gèle par la \emph{scrap iron clause} (un flux tronqué). Le capital ne finance pas une technique~; il finance un flux de revenus. Pas de flux, pas de déploiement.

Un troisième constat concerne la transformation du véhicule. La mise en série révèle que le véhicule financier n'est pas constant~: il change de forme avec chaque vague. Concession (câble), location-consommable (IBM), régulation \emph{utility} (Insull), brevet de principe (Bell). Chaque véhicule est une innovation institutionnelle qui résout le décalage temporel entre coûts et revenus d'une manière adaptée au type d'actif.

Ces trois constats sont des régularités, pas encore des mécanismes. Les formuler comme mécanismes exige de préciser les conditions sous lesquelles ils opèrent et les cadres théoriques qui en rendent compte.


\subsubsection{Trois mécanismes causaux}

\paragraph{Premier mécanisme --- L'investissement spéculatif fonctionnel}

Le télégraphe américain dans les années~1850 offre le cas le plus pur. \textcite{nonnenmacher_telegraph_2001} documente une concurrence anarchique~--- vingt compagnies, profits insignifiants, comptes mal tenus~--- suivie d'une consolidation sous Western Union en~1866 et d'un monopole profitable dans les années~1870. Le capital spéculatif avait financé un réseau que le capital prudent n'aurait pas construit. Le câble transatlantique suit la même séquence en accéléré~: 300\,000~livres investies, trois semaines de fonctionnement, faillite, recapitalisation, succès \parencite{gordon_thread_2002}. Le téléphone la reproduit sur trente ans~: rente de brevet à quarante-six pour cent \parencite{gabel_early_1969}, puis concurrence explosive (six mille compagnies en~1900), puis rachat par le syndicat Morgan-Vail \parencite{mueller_universal_1993}.

\markEVO \textcite{perez_technological_2002} propose le cadre qui rend compte de cette récurrence. Le capital financier joue un rôle fonctionnel dans la phase d'installation~: il prend des risques que le capital productif ne prendrait pas, parce que son horizon est plus court et sa capacité à sortir plus grande. La bulle n'est pas un dysfonctionnement~; elle est le mécanisme par lequel une société pré-finance une infrastructure dont les bénéfices sont collectifs et différés\footnote{La formulation est celle de Perez. Elle n'implique pas que toute bulle soit fonctionnelle~: Perez distingue explicitement les cas où l'infrastructure construite est réutilisable après la crise (chemins de fer, fibre optique) et ceux où elle ne l'est pas. La distinction dépend de la nature physique de l'actif, pas de l'intention de l'investisseur. \textcite{janeway_doing_2012} arrive à la même conclusion par l'expérience du praticien~: l'investissement spéculatif est une condition nécessaire pour construire l'infrastructure excédentaire qui servira de base à la vague suivante.}. Le mécanisme interne du glissement est celui de \textcite{minsky_stabilizing_1986}~: la logique \emph{hedge}~--- où les revenus couvrent le service de la dette~--- glisse vers une logique spéculative, puis éventuellement Ponzi, où seule la hausse des valorisations soutient le financement. \textcite{reinhart_this_2009} documentent le syndrome \og \emph{this time is different}\fg{} sur huit siècles de crises financières.

Les conditions d'activation sont précises~: rendements croissants de réseau (qui créent l'incitation au winner-take-all), coûts fixes élevés et irréversibles (qui rendent le capital \emph{sunk}), asymétrie informationnelle entre investisseurs et opérateurs (qui empêche l'évaluation correcte du rendement futur), et liquidité suffisante du marché financier pour alimenter la phase de frénésie. Quand l'une de ces conditions manque, le mécanisme ne s'enclenche pas~: le marché des calculatrices mécaniques, sans rendements de réseau, n'a connu ni bulle ni consolidation.

\paragraph{Hypothèse H-D8.1} \emph{Le déploiement d'une infrastructure informationnelle à rendements croissants attire un volume de capital qui excède ce que le calcul prudent ex ante justifierait, parce que les rendements anticipés de la phase d'installation enclenchent une dynamique auto-réalisatrice qui glisse de la position \emph{hedge} au sens de \textcite{minsky_stabilizing_1986} vers la position spéculative puis Ponzi, jusqu'à la rupture. La crise qui en résulte n'est pas un dysfonctionnement mais une étape du déploiement : elle transfère l'infrastructure à des opérateurs qui l'acquièrent à un coût inférieur au coût de construction, rendant le déploiement profitable ex post même quand il ne l'était pas ex ante \parencite{perez_technological_2002}.} Les conditions d'activation sont la conjonction de rendements croissants de réseau, de coûts fixes élevés et irréversibles, d'une asymétrie informationnelle entre investisseurs et opérateurs, et d'une liquidité suffisante du marché financier pour alimenter la phase de frénésie. Quand l'une de ces conditions manque~--- le marché des calculatrices mécaniques, sans rendements de réseau~--- le mécanisme ne s'enclenche pas.

\paragraph{Deuxième mécanisme --- L'adéquation du véhicule financier à la structure temporelle}

Le premier constat du matériau est que Babbage avait une technique fonctionnelle et un financement public unique de 17\,470~livres entre~1823 et~1842~; le câble avait une technique identique en~1857 et en~1866. Ce qui sépare l'échec du succès, c'est la présence d'un flux de revenus identifiable~--- l'arbitrage intercontinental~--- qui attire du capital neuf après la faillite. Le deuxième mécanisme porte sur cette variable~: non pas le volume du capital mais sa forme.

Toute infrastructure crée un décalage entre le moment où les coûts sont engagés (construction, ponctuelle et irréversible) et le moment où les revenus se matérialisent (exploitation, progressive et incertaine). Le véhicule financier résout ce décalage~: il convertit une dépense immédiate en un flux de revenus futur que l'investisseur accepte de financer. La forme du véhicule dépend de trois paramètres~: le ratio entre coûts fixes et coûts variables, l'horizon temporel des revenus, et la capacité de l'actif à servir de garantie (sa \emph{collatéralisabilité}).

\markINFO \textcite{hall_financing_2010} formalisent le problème pour l'innovation en général. \INFO{Le financement de l'innovation est contraint par trois propriétés~: l'incertitude sur les résultats, l'intangibilité de l'actif (qui ne peut pas servir de collatéral), et l'asymétrie informationnelle (l'inventeur en sait plus que l'investisseur).} \textcite{dell_ariccia_bank_2021} montrent que les banques réallouent leurs portefeuilles vers les actifs redéployables~--- ceux qu'elles peuvent saisir et revendre~--- quand les firmes basculent vers l'intangible. Le résultat est que la \emph{bancabilité} d'un actif dépend de sa redéployabilité, pas de sa productivité. Un data center est bancable parce que sa carcasse peut servir n'importe quel locataire~; un modèle de langage ne l'est pas. \textcite{esty_why_2004} formalise les conditions du project finance~: actifs hautement spécifiques, flux de revenus prévisibles sous contrat, levier d'endettement élevé. Ces trois cadres~--- Hall-Lerner sur le \emph{financing gap}, Dell'Ariccia sur la bancabilité, Esty sur le project finance~--- convergent vers un résultat commun~: la forme du capital informationnel (D3) détermine sa finançabilité.

Les six cas de l'analyse historique instancient cette logique. La \textbf{concession}~--- où l'État accorde un monopole temporaire contre l'obligation de construire~--- est le véhicule du câble transatlantique. Le \textbf{modèle de location avec consommable exclusif} est le véhicule d'IBM~: la tabulatrice à 40~dollars par mois et les cartes à 85~cents le millier transforment un \emph{capex} en \emph{opex} récurrent, et les cartes fournissent un flux de revenus indépendant du cycle d'investissement \parencite{cortada_before_1993}. La \textbf{régulation des \emph{utilities}} est le véhicule de l'électrification~: le régulateur garantit un rendement sur capital investi en échange d'une obligation de service, transformant un actif irréversible en quasi-bond \parencite{hughes_networks_1983}. Le \textbf{brevet de principe} est le véhicule du téléphone de Bell~: le juge Lowell qualifie la transmission de la parole de \og \emph{new art}\fg{} (1881), produisant une rente légale de quarante-six pour cent pendant dix-huit ans \parencite{gabel_early_1969}. Le \textbf{venture capital} est le véhicule de l'innovation intangible~: \textcite{hall_financing_2010} montrent que le VC résout l'asymétrie informationnelle par le contrôle actif, la sélection et le financement par étapes~; \textcite{lerner_venture_2020} montrent qu'il échoue pour la \emph{deep tech} (énergie, matériel physique) parce que les cycles sont trop longs et les besoins en capital trop élevés.

Ce qui donne au mécanisme sa force prédictive, c'est le test en négatif. Babbage n'avait aucun mécanisme de recapitalisation parce que la machine ne générait pas de flux de revenus. Le contraste avec le câble est structurel~: le câble avait une promesse de revenu qui attirait du capital nouveau après la faillite \parencite{gordon_thread_2002}~; Babbage n'avait qu'une promesse scientifique. À Londres, la \emph{scrap iron clause} de~1882 offrait aux investisseurs un rachat municipal à la valeur des composants séparés, sans prime pour la valeur du système en fonctionnement~: le véhicule ne détruisait pas l'infrastructure~; il détruisait l'incitation à la construire \parencite{hughes_networks_1983}. L'hypothèse abductive \#3 de l'analyse historique~--- la technique ne se déploie pas sans un véhicule de financement adapté~--- reçoit ici sa formulation causale~: l'adéquation ou l'inadéquation entre le véhicule et la structure des coûts est le mécanisme, pas le volume du capital disponible.

\paragraph{Hypothèse H-D8.2} \emph{Le déploiement d'une infrastructure informationnelle échoue quand le véhicule financier mobilisé est inadapté à la structure temporelle de ses coûts et de ses revenus, même si la technologie fonctionne et la demande existe. Cette adéquation se mesure sur trois paramètres : le ratio entre coûts fixes et coûts variables, l'horizon temporel des revenus, et la collatéralisabilité de l'actif~--- sa capacité à servir de garantie. L'incertitude, l'intangibilité et l'asymétrie informationnelle créent un \emph{financing gap} \parencite{hall_financing_2010} qui biaise l'allocation contre certaines configurations d'infrastructure.} La forme du capital informationnel (H-D2.1) détermine sa finançabilité : un actif tangible, collatéralisable et à revenus prévisibles (câble sous-marin) trouve un véhicule adapté ; un actif intangible non collatéralisable (algorithme) en trouve un autre ; un actif tangible à revenus imprévisibles (machine analytique de Babbage) n'en trouve aucun.

\paragraph{Troisième mécanisme --- La sélection financière entre infrastructures concurrentes}

Le classi-compteur de March remplissait la même fonction que la tabulatrice Hollerith pour le dépouillement statistique français~; il ne put se généraliser au-delà de quelques pays latins parce que son modèle financier~--- achat unique, pas de consommable, pas de revenus récurrents~--- ne permettait pas de financer une force de vente internationale comparable à celle de Watson \parencite{cortada_before_1993}. En France, le financement du réseau téléphonique par \og avances remboursables\fg{} des collectivités locales fragmenta le réseau selon la géographie administrative plutôt que selon la rationalité économique, produisant 14\,000~centraux pour 7\,000 en Allemagne \parencite{bertho_histoire_1984}. Dans ces deux cas, ce qui détermine le résultat n'est pas la qualité de l'infrastructure mais sa \emph{finançabilité}~--- l'adéquation entre les propriétés de l'actif et les critères du véhicule financier disponible.

\markNEO \textcite{philippon_has_2015} établit que le coût de l'intermédiation financière aux États-Unis est resté constant, autour de deux pour cent, depuis un siècle\footnote{\textcite{bazot_financial_2018} confirme le résultat pour l'Europe et montre que les banques utilisent l'ICT pour extraire de la rente de complexité plutôt que pour réduire les coûts d'intermédiation. \textcite{philippon_has_2015} lui-même documente un déclin récent, mais le puzzle reste largement non résolu.}. Cette constance agrégée masque une hétérogénéité considérable selon les types d'actifs. \textcite{haskel_capitalism_2018} montrent que les actifs intangibles ont quatre propriétés~--- \emph{scalability, sunkenness, spillovers, synergies}~--- qui rendent leur financement structurellement différent de celui des actifs tangibles. La \emph{scalability} attire le capital spéculatif~; la \emph{sunkenness} dissuade le capital prudent~; les \emph{spillovers} créent une divergence entre rendement privé et rendement social~; les \emph{synergies} concentrent les rendements sur les firmes qui possèdent déjà un portefeuille d'intangibles. Nous conjecturons que ces quatre propriétés produisent un biais de sélection financière~--- ce n'est pas l'argument de Haskel et Westlake, qui décrivent les propriétés sans en tirer cette conséquence pour l'allocation inter-infrastructures.

La conséquence est une sélection adverse du financement. L'infrastructure intangible à rendements croissants attire le capital-risque parce que sa \emph{scalability} permet des rendements multipliés par dix ou cent~; l'infrastructure tangible à coûts fixes élevés nécessite du financement de projet à rendements modestes et prévisibles. Les deux types d'infrastructure sont complémentaires dans le système technique, mais le système financier les finance séparément, avec des véhicules différents, à des horizons différents, et pour des investisseurs différents\footnote{La stagnation séculaire (Hansen~1939, Summers~2014) et la rareté des actifs sûrs \parencite{caballero_safe_2017} fournissent un arrière-plan macroéconomique cohérent avec H-D8.3~: l'insuffisance d'actifs sûrs déprime leur rendement et pousse le capital vers les actifs risqués à rendement élevé. L'infrastructure tangible, qui pourrait servir d'actif sûr, n'est pas \emph{packagée} pour ce marché~; l'infrastructure intangible capte le capital excédentaire. Le lien entre stagnation séculaire et biais de financement de l'infrastructure numérique reste conjectural et n'a pas, à notre connaissance, été établi explicitement dans la littérature.}.

Les conditions d'activation sont~: existence simultanée de deux types d'infrastructure en compétition pour le capital, véhicules financiers spécialisés par type d'actif, et absence de mécanisme institutionnel qui internalise la complémentarité entre les deux. Les câbles transpacifiques de~1902 illustrent le contraire~: quand il n'y a qu'un type d'infrastructure en compétition avec lui-même, le biais ne joue pas~--- et le filtre est la demande réelle (\textcite{muller_telegraph_2015} montrent que les câbles pacifiques échouèrent parce que le trafic n'existait pas).

\paragraph{Hypothèse H-D8.3} \emph{Quand plusieurs formes d'infrastructure informationnelle sont en compétition pour le même capital, l'allocation suit la \emph{finançabilité}~--- l'adéquation entre les propriétés de l'actif et les critères du véhicule financier disponible~--- et non la valeur sociale. La spécialisation des véhicules par type d'actif (\emph{venture capital} pour l'intangible scalable, \emph{project finance} pour le tangible à revenus prévisibles) produit un biais systématique en faveur des infrastructures à forte \emph{scalability} et au détriment des infrastructures tangibles à rendements modestes et prévisibles, même quand les deux sont complémentaires dans le système technique.} Les conditions d'activation sont la coexistence de deux types d'infrastructure en compétition pour le capital, la spécialisation des véhicules financiers par type d'actif, et l'absence de mécanisme institutionnel qui internalise la complémentarité. Le cas négatif est le câble transpacifique de 1902 : un seul type d'infrastructure en compétition avec lui-même, le biais ne joue pas~--- l'échec vient de la demande, pas du financement \parencite{muller_telegraph_2015}.


\subsubsection{Test contemporain}

\paragraph{Évidence historique}

Les cas de l'analyse historique instancient les trois mécanismes selon des combinaisons distinctes, ce qui permet de vérifier les conditions d'activation plutôt que de les illustrer.

\emph{Le câble transatlantique (M1~+~M2).} Le câble de Cyrus Field suit le pattern complet du premier mécanisme. Le financement initial mobilisa des capitaux britanniques considérables sur une promesse de rente informationnelle~: celui qui contrôlerait le câble contrôlerait l'arbitrage coton-blé-titres entre les deux continents. Le premier câble (1857--1858) coûta environ 300\,000~livres à l'\emph{Atlantic Telegraph Company}~; il ne fonctionna que trois semaines et la compagnie fit faillite. La recapitalisation de~1864, avec les fonds de Pender et Brassey, leva le capital nécessaire à un câble réussi en~1866 pour 600\,000 à 700\,000~livres \parencite{gordon_thread_2002}. Le deuxième mécanisme opère en parallèle~: le véhicule financier~--- la concession~--- était adapté à un actif dont les coûts fixes étaient énormes et les revenus concentrés sur un corridor commercial dense. Les câbles transpacifiques de~1902--1904 fournissent le contre-test~: le véhicule était le même, mais le corridor n'existait pas~; les comptes ne rentraient pas \parencite{muller_telegraph_2015}.

\emph{Babbage~: l'anti-pattern (M2 négatif).} Le financement public (17\,470~livres entre~1823 et~1842) ne comportait aucun mécanisme de recapitalisation parce que la machine ne générait pas de revenus. Le \emph{Difference Engine} produisait des tables dont la demande était étroite~--- un exemplaire unique \parencite{swade_difference_2001}. Le contraste avec le câble est structurel~: le câble avait une promesse de revenus qui mobilisait une recapitalisation~; Babbage n'avait qu'une promesse scientifique. La \emph{scrap iron clause} londonienne de~1882 offre un deuxième anti-pattern, symétrique~: la demande existait, mais le véhicule détruisait l'incitation en amputant la valeur de revente de l'actif \parencite{hughes_networks_1983}.

\emph{IBM~: le véhicule adapté (M2 positif, M3).} Le modèle de location avec consommable exclusif est le cas le plus pur du deuxième mécanisme. Les revenus d'IBM passèrent de 2,5~millions de dollars en~1919 à 45,3~millions en~1940, et la compagnie paya des dividendes sans interruption pendant la Grande Dépression \parencite{cortada_before_1993}. Le troisième mécanisme opère aussi~: le classi-compteur de March, qui remplissait la même fonction, ne put se généraliser parce que son modèle financier ne permettait pas de financer une force de vente internationale.

\emph{Le téléphone~: le cycle complet (M1~+~M2~+~M3).} Le cas du téléphone combine les trois mécanismes en séquence. Le monopole de brevet de Bell (1876--1894) instancie un véhicule fondé sur la rente légale à quarante-six pour cent \parencite{gabel_early_1969}. L'expiration ouvrit la concurrence~: plus de six mille compagnies en~1900~; la crise des profits conduisit le syndicat Morgan-Vail à racheter et reconsolider le réseau à partir de~1907 \parencite{mueller_universal_1993}. Le \emph{Kingsbury Commitment} (1913) institutionnalisa le monopole sous couvert de régulation~--- un véhicule financier déguisé en accord antitrust. En France, le financement par \og avances remboursables\fg{} fragmenta le réseau~: 14\,000~centraux pour 7\,000 en Allemagne \parencite{bertho_histoire_1984}. Le véhicule détermina la forme du réseau.


\paragraph{Évidence contemporaine}

Le marqueur financier le plus net de la période contemporaine est le passage du venture capital au project finance comme véhicule dominant de l'investissement en infrastructure informationnelle. Pendant trente ans (1990--2020), le secteur numérique a été financé par le VC et par l'autofinancement des grandes firmes technologiques. Les dépenses en capital des quatre principaux hyperscalers (Amazon, Microsoft, Alphabet, Meta) atteignent environ 225~milliards de dollars en~2024 et les guidances pour~2025 dépassent 320~milliards\footnote{Sources~: SEC filings 10-K et 10-Q~; agrégation par \emph{Platformonomics} (Q3~2025 Scoreboard). Amazon~: 77~Md\$ (2024), 100--105~Md\$ (guidance 2025)~; Microsoft~: 55,7~Md\$ (FY2024), 80~Md\$ (guidance FY2025)~; Alphabet~: 52~Md\$ (2024), 75~Md\$ (guidance 2025)~; Meta~: 40~Md\$ (2024), 64--72~Md\$ (guidance 2025). L'autofinancement domine~: le cash-flow opérationnel de Microsoft (${\sim}$110~Md\$ en FY2024) couvre largement le capex. L'émission de dette croît mais reste secondaire.}. L'émergence des data centers à l'échelle du gigawatt et des contrats d'achat d'énergie de long terme marque un basculement~: l'infrastructure informationnelle redevient tangible, et avec elle le véhicule financier bascule.

Ce basculement instancie le deuxième mécanisme. \textcite{lerner_venture_2020} montrent que le VC est structurellement inadapté au financement d'actifs physiques à cycles longs et besoins en capital élevés. Le capital-risque, dont l'horizon est de cinq à sept ans et dont la logique est celle de la sortie par cession, ne convient pas à un actif dont le coût de fonctionnement est permanent. Le passage au project finance est la réponse~: \textcite{esty_why_2004} formalise les conditions~--- actifs spécifiques, flux de revenus prévisibles sous contrat, levier d'endettement élevé. Le \emph{Global AI Infrastructure Investment Partnership} (GAIIP), lancé en septembre~2024 par BlackRock (GIP), Microsoft et MGX, incarne ce basculement~: 30~milliards de dollars en equity pour mobiliser jusqu'à 100~milliards avec de la dette\footnote{Le GAIIP est le premier véhicule de project finance dédié à l'infrastructure IA. Sa structure~--- equity d'un fonds d'infrastructure + garantie corporate d'un hyperscaler + levier bancaire~--- est un hybride qui comble le gap entre le VC (horizon trop court) et le project finance classique (actif trop nouveau). Historiquement, le syndicat Morgan-Vail de~1907 remplissait une fonction comparable~: du capital bancaire qui rachète une infrastructure fragmentée pour la consolider.}. Les spreads de crédit des \emph{asset-backed securities} adossées aux data centers se situent à 150--200 points de base au-dessus du taux sans risque, un niveau comparable aux infrastructures \emph{core} (autoroutes, aéroports)\footnote{Données de pricing~: L\&G Asset Management, \og Securitized Credit~: Financing Data Centers\fg{} (2026). La compression récente des spreads à ${\sim}$160~bps signale que le marché traite les data centers comme une classe d'actifs d'infrastructure, pas comme un risque technologique.}. Les contrats d'achat d'énergie (PPAs) jouent un rôle analogue aux concessions du dix-neuvième siècle~: ils garantissent un flux de revenus fixe sur dix à vingt-cinq ans, ce qui permet aux banques de prêter à des ratios de levier de 70 à 80\,\%.

H-D8.3 prédit que ce basculement crée un biais d'allocation. Le venture capital continue de financer les couches logicielles (modèles, applications, interfaces) tandis que le project finance finance les couches physiques (data centers, réseaux, énergie). Les deux couches sont complémentaires dans le système technique mais financées par des véhicules séparés, pour des investisseurs séparés, à des horizons séparés. Le résultat est un risque structurel de sous-investissement dans la couche physique relativement à la couche logicielle, parce que le rendement de la couche logicielle est plus visible, plus rapide, et plus scalable que celui de la couche physique. La portée de cette prédiction reste conjecturale~: elle dépend de la persistance du modèle de bilan des hyperscalers et de la stabilité du cadre réglementaire qui distingue VC et project finance.

L'analogie avec le schéma de \textcite{perez_technological_2002} (irruption $\rightarrow$ frénésie $\rightarrow$ synergie $\rightarrow$ maturité) appliquée au cycle d'investissement IA mérite prudence. La phase actuelle (2023--2026) présente certains symptômes de la frénésie~--- investissements massifs, valorisations déconnectées des revenus, syndrome \og \emph{this time is different}\fg{}. Mais une objection structurelle s'oppose à l'identification complète~: les hyperscalers qui financent l'infrastructure disposent de cash-flows opérationnels considérables et financent une part substantielle de l'investissement sur fonds propres, non par endettement spéculatif. Le risque principal n'est donc pas la faillite minskyenne mais l'allocation massive de capital dans une infrastructure dont le coût de fonctionnement est permanent et dont les revenus restent partiellement non garantis. La question de savoir si le cycle actuel suivra le pattern Perez complet~--- crise financière suivie d'un déploiement productif~--- ou s'il représente un régime nouveau~--- investissement autofinancé dont le risque est l'allocation sous-optimale plutôt que la faillite~--- reste ouverte\footnote{L'État est actif dans le cycle contemporain. Le CHIPS and Science Act (2022, 52,7~milliards de dollars) a déclenché environ 540~milliards de dollars d'investissements privés annoncés, soit un ratio privé/public d'environ 10:1. L'European Chips Act (2023, 43~milliards d'euros) reproduit le pattern. Le rôle de l'État comme preneur de risque de premier rang dans l'infrastructure informationnelle~--- documenté par \textcite{flamm_creating_1988} pour la R\&D informatique américaine des années~1950 (environ deux tiers du financement)~--- se confirme. Les mécanismes de retour vers le financeur public ne sont pas intégrés dans les conditions de ces subventions~: le pattern historique~--- socialisation du risque, privatisation du rendement \parencite{mazzucato_entrepreneurial_2013}~--- se reproduit.}.


\subsubsection{Fait stylisé}

L'infrastructure informationnelle ne se déploie pas sans un véhicule de financement adapté à sa structure de coûts. Le capital ne finance pas une technique~; il finance un flux de revenus~--- cette formulation est notre synthèse à partir du matériau historique et des cadres de \textcite{perez_technological_2002}, \textcite{hall_financing_2010} et \textcite{esty_why_2004}. La forme du véhicule --- concession, location avec consommable, régulation des utilities, brevet de principe, venture capital, project finance --- est déterminée par trois paramètres~: le ratio coûts fixes / coûts variables, l'horizon temporel des revenus, et la collatéralisabilité de l'actif (H-D8.2). Quand le véhicule est inadapté (Babbage~: pas de flux de revenus~; \emph{scrap iron clause}~: incitation détruite~; câbles transpacifiques~: pas de demande commerciale), la technologie échoue commercialement même quand elle fonctionne techniquement. Quand il est adapté (IBM~: location-consommable~; Insull~: \emph{utility regulation}~; câble atlantique~: concession sur corridor commercial dense), le déploiement procède. La séquence investissement spéculatif, crise, recapitalisation n'est pas un dysfonctionnement mais un mécanisme récurrent de transfert d'infrastructure (H-D8.1), documenté pour le télégraphe, les chemins de fer et les dotcoms \parencite{perez_technological_2002, nonnenmacher_telegraph_2001, reinhart_this_2009}. Quand des infrastructures concurrentes se disputent le même capital, l'allocation suit la finançabilité, pas la valeur sociale (H-D8.3). Le passage contemporain du venture capital au project finance~--- dont le GAIIP de septembre~2024 est le marqueur institutionnel~--- est le signe financier du retour de la tangibilité dans l'infrastructure informationnelle. Les spreads de crédit des data centers convergent vers ceux des autoroutes et des aéroports~: le marché financier a reclassé l'infrastructure informationnelle comme infrastructure lourde avant que la théorie économique ne le fasse.

Le retour de la tangibilité rend pertinente la distinction entre un capital physique numérique dont le financement obéit à des contraintes de bilan d'hyperscaler~--- horizon long, actifs collatéralisables, flux contractuels~--- et un capital intangible financé par VC. La structure de financement n'est pas représentée explicitement dans les cadres CGE standards~; elle peut entrer comme contrainte de calibration sur les durées de vie et les paramètres putty-clay du capital informationnel, ce qui constitue un paramétrage de scénario examinable dans la modélisation. Le biais d'allocation entre couches (H-D8.3) reste comme contrainte interprétative à expliciter dans la lecture des trajectoires.


% =============================================================================
